\subsection{全supportと正常一次hdetによる20有限ラベルへの削減}
\label{subsec:cubic-four-simple-support-sieve}
\begin{lemma}[連結二三次型の共通因子]
\label{lem:cubic-linked-binary-factor}
相異なる一次元 $u,v,z$ について、$uz$ が $zu$ の非零scalar倍であり、
二次関係が一本
\[
 a\,vz+b\,zv+c\,u^2=0
\]
存在するとする。整域性から $(a,b)\ne(0,0)$ である。
三次 $G$ の全語の指数が $u^2z,vz^2$ のいずれかなら、$G$ は共通一次因子 $z$ を持つ。
\end{lemma}
\begin{proof}
$b\ne0$ なら $z^2v=z(zv)$ を $zvz,zu^2$ に書き換える。
$u,z$ のskew可換性により、二型の全語は右因子 $z$ を持つ。
$b=0$ なら $a\ne0$、$vz$ を $u^2$ に書き換えれば全語は左因子 $z$ を持つ。
$c=0$ の場合も同じ議論である。
\end{proof}

\begin{proposition}[正常化scalarの整数左核障害]
\label{prop:cubic-support-normalizer-leftkernel}
全二次元pair関係から、いくつかの生成元 $x_z$ が全生成元と個別にskew可換であることが強制されているとする。
$x_ix_j=q_{ij}x_jx_i$ の向きを $i<j$ に固定し、正常元に接する強制pairの $q_{ij}$ を未知非零scalarとする。
実際のtopの非零文字数supportを $S$、正常三次 $G$ の共通固有liftの文字数supportを $T$ とする。
各正常生成元 $x_z$ と $w\in S\cup T$ に対して
\[
 \prod_{i>z}q_{zi}^{w_i}\prod_{i<z}q_{iz}^{-w_i}=\lambda_z
\]
が必要である。
固有値を $\lambda_z=\zeta_N^{a_z}$ と書き、左辺の指数行列を $A$、右辺の整数指数列を $b$ とする。
整数行 $\ell$ が $\ell A=0$ かつ $\ell b\not\equiv0\pmod N$ を満たせば、この全support chartは空である。
\end{proposition}
\begin{proof}
強制pairの各両係数が非零なので、正常化自己同型 $\phi_z$ は対角で、
$\phi_z(x_z)=x_z$、他文字のscalarは表示した $q_{ij}$ 又は逆数である。
補題\ref{lem:cubic-normal-linear-hdet}より、そのhdetは $\lambda_z$。
全非零top語は同じhdetで変換されるから $w\in S$ の式を得る。
また $Gx_z=\lambda_z^{-1}x_zG$ と整域性から $\phi_z(G)=\lambda_zG$。
全 $\phi_z$ と中山は対角なので互いに可換であり、tensorの各次数部分とその商で同時半単純である。
従って $G$ は全 $\phi_z$ と中山の共通固有空間からliftでき、$w\in T$ の各語も同じ式を満たす。
恣意的なliftの各単項式に固有値条件を課した議論ではない。

表示したmonomial方程式を整数 $\ell$ で掛け合わせると、左辺は1、右辺は $\zeta_N^{\ell b}$。
これが1でないことが矛盾である。
根の数値近似も、未知 $q_{ij}$ の有限位数性も仮定していない。
\end{proof}

\begin{theorem}[20有限ラベルと全1306必要support片]
\label{thm:cubic-four-simple-twenty-support-atlas}
全四単純H3スペクトルの必要被覆は次の20 Galoisラベル軌道、45根集合へ縮む。
従って中山の位数は $5,6,7,8,9,12$ のいずれかである。
各根集合上の全候補は、公開証明書の1306個の必要support片で覆われる。
各片は二次関係character、top文字数support、三次文字数supportで指定される。
この片の非空性、AS4性又はH3許容性は主張しない。
\end{theorem}
\begin{proof}
前定理の39ラベルそれぞれで、残る全関係characterを固定する。
元topの可能な文字数の有限集合から、非空な全部分集合 $S$ を取る。
全変数を被覆しないもの、multilinearだけのものは除く。
関係rank零のblockに非零成分を持つ巡回軌道は、その係数が零となるので除く。
次に一monomial列のdomain pruningを安定まで反復する。
$S$ 中の文字数型が全て消えた場合、その $S$ は実際のsupportでないので除く。

残る各blockの非零位置を二部graphとする。
行列rankはこのgraphの最大matching以下なので、指定 $r_s$ がこの上界を超えるchartは空である。
ここでmatchingをrank十分性と取り違えてはいない。
像の可能な行supportが一つのpair $ij,ji$ だけで、$r_s=1$ なら、
整域性によりそのpairは非零skew関係として強制される。
この全強制pairから、全正常生成元とそれらに接する未知scalarを決める。
六pair全部が強制skewなら既分類に帰着して除く。

次に固定三次の全可能文字数から全非空部分集合 $T$ を取る。
前節の共通skew正常一次因子、純二元AS2因子、及び補題\ref{lem:cubic-linked-binary-factor}の必要因子があれば除く。
残りについて前命題の整数行列を作り、Smith標準形を計算する。
左・右の変換行列が整数単元であることと $UAV=D$ の全恒等式を確認し、
零対角行に対応する $U$ の各行を整数左核として使う。
剰余が非零なら、その行 $\ell$ と $\ell b\pmod N$ を証明書へ保存する。
公開checkerはさらに全保存行について、Smith実装に依存せず
$\ell A=0$ と剰余非零を直接再検算する。

この操作は全 $S,T$ を走査するので、各実際のH3対象は必ず一つの残る片に入る。
残るcharacterごとの走査も全て行い、一般点から特殊supportを消してはいない。
39ラベル中19ラベルでは一つの片も残らず、20ラベルに全1306片が残る。
下表は $N$、指数代表、残る必要support片の数である。
Galois共役を展開すると45根集合となる。
\end{proof}

\begin{center}
\begin{longtable}{r l r}
\toprule
$N$&指数代表&必要support片数\\\midrule
\endfirsthead
\toprule $N$&指数代表&必要support片数\\\midrule
\endhead
5&$(1,2,3,4)$&633\\
6&$(0,1,2,3)$&50\\
6&$(0,1,3,5)$&64\\
6&$(0,2,3,4)$&424\\
6&$(1,2,4,5)$&36\\
7&$(0,1,2,4)$&1\\
7&$(1,2,5,6)$&15\\
8&$(0,1,5,6)$&3\\
8&$(1,2,3,6)$&1\\
8&$(1,2,4,5)$&2\\
8&$(1,2,6,7)$&3\\
9&$(0,1,4,7)$&14\\
9&$(1,3,4,7)$&7\\
9&$(1,3,6,8)$&15\\
9&$(1,4,6,7)$&19\\
12&$(0,1,4,10)$&8\\
12&$(1,3,4,7)$&2\\
12&$(1,3,7,10)$&2\\
12&$(1,4,6,10)$&4\\
12&$(3,4,8,9)$&3\\
\bottomrule
\end{longtable}
\end{center}

公開再現器は \texttt{verify\_cubic\_four\_simple\_support\_sieve.py}。
全support走査と整数Smithは \texttt{cubic\_four\_simple\_support\_sieve.py} にある。
各片は、証明書に記した文字数型配列に対するbitmaskと関係character番号で保持する。
適合した片を実際の正常incidenceと呼ばず、次に元middle rank式と三次正常式を判定する。
全CY・成分・境界・intrinsic群oidと相対降下も未完成である。
